\chapter{Алгебра Стандартной модели, семейный коммутант и предел обычных одноформ}

Предыдущая глава восстановила единый $300$-мерный носитель, но оставила
открытым способ совместить семейную цепь с наблюдаемой координатной алгеброй.
Теперь проверим наиболее экономное чтение старых результатов:
\begin{equation}
 \mathcal A_{\mathrm{coord}}=\mathcal A_{\mathrm{SM}}
 =\mathbb C\oplus\mathbb H\oplus M_3(\mathbb C),
 \label{eq:v5-sm-coordinate-algebra}
\end{equation}
а аффинная семейная цепь действует не как второй координатный множитель, а
как градуированное соответствие в коммутанте представления
$\pi(\mathcal A_{\mathrm{SM}})$.

Такое разделение согласуется с диаграммами Краевского: конечная алгебра
задаёт левое и правое действие на бимодуле, а число поколений может входить
как кратность представления. Оно также близко к языку неограниченных
соответствий Каспарова, где морфизм между спектральными геометриями несёт
собственный оператор и связность. Однако литературная аналогия ещё не
доказывает, что семейная цепь возникает из обычных внутренних флуктуаций.

\section{Коммутантное размещение}

Пусть $K_{\mathrm{fam}}$ --- частичный десятимерный носитель цепи
$4\to3\to3$, а $H_{\mathrm{SM}}$ --- пятнадцатимерный наблюдаемый пакет.
На
\begin{equation}
 H_p=K_{\mathrm{fam}}\otimes H_{\mathrm{SM}}
\end{equation}
зададим
\begin{equation}
 \pi(a)=I_{10}\otimes\pi_{\mathrm{SM}}(a),
 \qquad
 D=D_{\mathrm{fam}}\otimes I_{15}
   +I_{10}\otimes D_{\mathrm{SM}}.
 \label{eq:v5-sm-family-product}
\end{equation}
После KO6-удвоения размерность равна
\begin{equation}
 2\cdot10\cdot15=300.
\end{equation}
Семейный оператор коммутирует со всей координатной алгеброй:
\begin{equation}
 [D_{\mathrm{fam}}\otimes I_{15},\pi(a)]=0.
 \label{eq:v5-family-in-commutant}
\end{equation}
Поэтому калибровочная алгебра не получает десяти независимых копий и после
условия унимодулярности остаётся
\begin{equation}
 \mathfrak{su}(3)\oplus\mathfrak{su}(2)\oplus\mathfrak u(1).
\end{equation}
Это решает кинематическую часть задачи: семейный $M_3(\mathbb R)$ не
перемножается с цветовым $M_3(\mathbb C)$.

\section{Теорема о невидимости семейной цепи}

\begin{theorem}[Невидимость коммутантного оператора для обычных одноформ]
Для представления и оператора \eqref{eq:v5-sm-family-product} выполнено
\begin{equation}
 [D,\pi(a)]
 =I_{10}\otimes[D_{\mathrm{SM}},\pi_{\mathrm{SM}}(a)].
 \label{eq:v5-one-form-blindness}
\end{equation}
Следовательно, пространство обычных одноформ
\begin{equation}
 \Omega_D^1(\mathcal A_{\mathrm{SM}})
 =\left\{\sum_j\pi(a_j)[D,\pi(b_j)]\right\}
\end{equation}
не зависит от $D_{\mathrm{fam}}$. Ни изменение семейной цепи, ни её
аффинный вакуум не порождают через это исчисление семейного скалярного поля.
\end{theorem}

\begin{proof}
Первое слагаемое $D$ коммутирует с $I_{10}\otimes\pi_{\mathrm{SM}}(a)$,
а коммутатор второго слагаемого равен правой части
\eqref{eq:v5-one-form-blindness}. Умножение слева на элементы алгебры
сохраняет общий множитель $I_{10}$. Поэтому замена
$D_{\mathrm{fam}}\mapsto D'_{\mathrm{fam}}$ не меняет ни одного генератора
одноформ.
\end{proof}

Машинный контроль на верном пятнадцатимерном полупростом представлении дал
нулевые дефекты всех трёх тождеств. Размер комплексной линейной оболочки
контрольных одноформ равен $51$ как до, так и после десятикратного
коммутантного размещения. Само число $51$ зависит от контрольного
представления; его неизменность --- нет, поскольку следует из теоремы.

\section{Трилемма объединения алгебр}

Остаются три очевидных способа соединения наблюдаемого и семейного слоёв.

\subsection{Прямая сумма}

Для
\begin{equation}
 \mathcal A_{\mathrm{SM}}\oplus\mathcal A_{\mathrm{fam}}
\end{equation}
центральные единицы двух слагаемых ортогональны:
$e_{\mathrm{SM}}e_{\mathrm{fam}}=0$. Неприводимый модуль принадлежит одному
центральному сектору и не может быть одновременно нетривиально заряжен по
обоим слагаемым. Связующий сектор тогда вновь требует дополнительного
бимодуля и внешней нормировки.

\subsection{Тензорное произведение}

Наивное тензорное объединение создаёт матричный блок
\begin{equation}
 M_3(\mathbb R)\otimes_{\mathbb R}M_3(\mathbb C)
 \simeq M_9(\mathbb C).
\end{equation}
На уровне вещественных алгебр Ли отдельные направления
$\mathfrak{so}(3)\oplus\mathfrak u(3)$ имеют размерность $12$, тогда как
$\mathfrak u(9)$ имеет размерность $81$. Возникают $69$ лишних генераторов.
Это не тонкая феноменологическая поправка, а смена калибровочной теории.

\subsection{Коммутант}

Коммутант сохраняет правильную наблюдаемую калибровочную алгебру, но по
теореме \eqref{eq:v5-one-form-blindness} обычное исчисление не видит
семейный дифференциал. Поэтому этот вариант проходит кинематику и
одновременно проваливает происхождение семейной связности внутри одной
обычной спектральной тройки.

\begin{nogocriterion}[Обычное коммутантное замыкание]
Если семейный оператор целиком принадлежит коммутанту координатной алгебры,
то он не может быть источником её внутренних флуктуаций. Добавлять семейный
потенциал после этого отдельным членом запрещено правилом отсутствия
последующей подгонки.
\end{nogocriterion}

\section{Что действительно остаётся открытым}

Трилемма не закрывает программу. Она уточняет математический тип следующего
объекта. Семейная цепь должна быть не вторым координатным множителем и не
невидимой добавкой к $D$, а самостоятельным градуированным соответствием
над $\mathcal A_{\mathrm{SM}}$ со связностью
\begin{equation}
 \nabla_{\mathrm{fam}}:
 K_{\mathrm{fam}}\longrightarrow
 K_{\mathrm{fam}}\widehat\otimes_{\mathcal A_{\mathrm{SM}}}
 \Omega^1(\mathcal A_{\mathrm{SM}}),
\end{equation}
либо компонентой градуированной суперсвязности. Тогда ориентированный
дифференциал $d_{\mathrm{fam}}$, высота и отображение момента являются
данными морфизма, а не искусственно вставленным знаком в спектральный след.

Это пока не решение, а строго выделенная ветвь. Следующий гейт обязан
проверить одновременно:
\begin{enumerate}
 \item градуированное правило Лейбница;
 \item совместимость с $J$, градуировкой и KO6;
 \item ковариантность относительно калибровочной группы Стандартной модели;
 \item происхождение квадрата кривизны из одного положительного следа
 $\tau_{300}$;
 \item кинетическую нормировку семейных скаляров из того же действия;
 \item полный комплекс БВ/БРСТ без дополнительных весов.
\end{enumerate}

\begin{protocol}
\begin{enumerate}
 \item алгебра Стандартной модели как единственная координатная:
 \textbf{сохранена};
 \item размерность полного KO6-носителя:
 \textbf{$300$};
 \item отсутствие размножения калибровочной алгебры в коммутанте:
 \textbf{доказано};
 \item происхождение семейной цепи из обычных одноформ:
 \textbf{опровергнуто};
 \item прямая сумма:
 \textbf{закрыта центральным разложением};
 \item наивное тензорное произведение:
 \textbf{закрыто калибровочным расширением};
 \item градуированное соответствие со связностью:
 \textbf{остаётся открытым};
 \item физическое замыкание:
 \textbf{не достигнуто}.
\end{enumerate}
\end{protocol}

Машинная проверка и сертификат сохранены в
\path{s2t_v5_sm_family_commutant_calculus_gate.py} и
\path{s2t_v5_sm_family_commutant_calculus_gate_results.json}.